% GENERATED FILE. Edit canonical lessons, then run npm run generate:books.
% canonical-source-hash: fnv1a64:e9d01cf807b208a0
% canonical-lessons: TA-C406-eriporul, TA-C406-petrol, TA-C406-tayar, TA-C406-tir, TA-C406-tey

\chapter{Fuel and Tyres}
\label{ch:ta-fuel-and-tyres}
\hlchaptermodality{406}

\begin{chapteropening}
\textbf{By the end of this chapter:} \emph{I can talk about fuel, petrol and tyres, and running out and wearing down.}

Complete the last lesson of chapter 406: I can talk about fuel, petrol and tyres, and running out and wearing down.
\end{chapteropening}

\section[\texorpdfstring{\ta{எரிபொருள்} (eriporuḷ)}{எரிபொருள் (eriporuḷ)}]{\ta{எரிபொருள்} (eriporuḷ) — fuel}

\label{lesson:TA-C406-eriporul}



\begin{quote}
Before the new one: say the Tamil for a lounge, then the Tamil for an interval.
\end{quote}

\subsection*{You'll want to know: \ta{எரிபொருள்}}
\textbf{\ta{எரிபொருள்}} — \emph{eriporuḷ} — \textquotedblleft{}fuel\textquotedblright{}.

\begin{quote}
\ta{எங்கள்} \ta{காரில்} \ta{எரிபொருள்} \ta{குறைவாக} \ta{இருக்கிறது}.
\end{quote}

\begin{quote}
\emph{eṅkaḷ kāril eriporuḷ kuṟaivāka irukkiṟatu.}
\end{quote}

\begin{quote}
Our car is low on fuel.
\end{quote}

\ta{எரி}, to burn, + \ta{பொருள்}, a thing: a thing that burns.

\begin{cousinweb}
\textbf{\ta{எரிபொருள்}} (\emph{eriporuḷ}) is built from \ta{எரி}, \textquotedblleft{}to burn\textquotedblright{} (DEDR 811), and \ta{பொருள்}, \textquotedblleft{}a thing\textquotedblright{} (DEDR 4544).
\end{cousinweb}

\begin{grammarlens}[title={What you've built}]
One more word for getting around.
\end{grammarlens}

\subsection*{Your turn}
\begin{itemize}
\raggedright
  \item \emph{Say it:} \emph{eriporuḷ}
  \item \emph{Say it:} \emph{eriporuḷ}, once more
  \item \emph{Say it:} \emph{eṅkaḷ kāril eriporuḷ kuṟaivāka irukkiṟatu.} — Our car is low on fuel.
  \item \emph{Recall:} say the Tamil for a lounge, then the Tamil for an interval, then say \emph{eriporuḷ} again
\end{itemize}

\begin{tcolorbox}[breakable,colback=teal!4,colframe=teal!35!black,title={Before you move on}]
What does \ta{எரிபொருள்} mean? (\textquotedblleft{}Fuel\textquotedblright{}.) Say it once more.
\end{tcolorbox}
\section[\texorpdfstring{\ta{பெட்ரோல்} (peṭrōl)}{பெட்ரோல் (peṭrōl)}]{\ta{பெட்ரோல்} (peṭrōl) — petrol}

\label{lesson:TA-C406-petrol}



\begin{quote}
Before the new one: say the Tamil for an interval, then the Tamil for fuel.
\end{quote}

\subsection*{You'll want to know: \ta{பெட்ரோல்}}
\textbf{\ta{பெட்ரோல்}} — \emph{peṭrōl} — \textquotedblleft{}petrol\textquotedblright{}.

\begin{quote}
\ta{பெட்ரோல்} \ta{விலை} \ta{மீண்டும்} \ta{ஏறியது}.
\end{quote}

\begin{quote}
\emph{peṭrōl vilai mīṇṭum ēṟiyatu.}
\end{quote}

\begin{quote}
The price of petrol went up again.
\end{quote}

Prices climb in Tamil: \ta{விலை} \ta{ஏறியது}, the price went up.

\begin{cousinweb}
\textbf{\ta{பெட்ரோல்}} (\emph{peṭrōl}) is a loan from English \emph{petrol}.
\end{cousinweb}

\begin{grammarlens}[title={What you've built}]
One more word for getting around.
\end{grammarlens}

\subsection*{Your turn}
\begin{itemize}
\raggedright
  \item \emph{Say it:} \emph{peṭrōl}
  \item \emph{Say it:} \emph{peṭrōl}, once more
  \item \emph{Say it:} \emph{peṭrōl vilai mīṇṭum ēṟiyatu.} — The price of petrol went up again.
  \item \emph{Recall:} say the Tamil for an interval, then the Tamil for fuel, then say \emph{peṭrōl} again
\end{itemize}

\begin{tcolorbox}[breakable,colback=teal!4,colframe=teal!35!black,title={Before you move on}]
What does \ta{பெட்ரோல்} mean? (\textquotedblleft{}Petrol\textquotedblright{}.) Say it once more.
\end{tcolorbox}
\section[\texorpdfstring{\ta{டயர்} (ṭayar)}{டயர் (ṭayar)}]{\ta{டயர்} (ṭayar) — a tyre}

\label{lesson:TA-C406-tayar}



\begin{quote}
Before the new one: say the Tamil for fuel, then the Tamil for petrol.
\end{quote}

\subsection*{You'll want to know: \ta{டயர்}}
\textbf{\ta{டயர்}} — \emph{ṭayar} — \textquotedblleft{}a tyre\textquotedblright{}.

\begin{quote}
\ta{வழியில்} \ta{காரின்} \ta{டயர்} \ta{வெடித்தது}.
\end{quote}

\begin{quote}
\emph{vaḻiyil kāriṉ ṭayar veṭittatu.}
\end{quote}

\begin{quote}
On the way the car's tyre burst.
\end{quote}

The plural is \ta{டயர்கள்}.

\begin{cousinweb}
\textbf{\ta{டயர்}} (\emph{ṭayar}) is a loan from English \emph{tyre}.
\end{cousinweb}

\begin{grammarlens}[title={What you've built}]
One more word for getting around.
\end{grammarlens}

\subsection*{Your turn}
\begin{itemize}
\raggedright
  \item \emph{Say it:} \emph{ṭayar}
  \item \emph{Say it:} \emph{ṭayar}, once more
  \item \emph{Say it:} \emph{vaḻiyil kāriṉ ṭayar veṭittatu.} — On the way the car's tyre burst.
  \item \emph{Recall:} say the Tamil for fuel, then the Tamil for petrol, then say \emph{ṭayar} again
\end{itemize}

\begin{tcolorbox}[breakable,colback=teal!4,colframe=teal!35!black,title={Before you move on}]
What does \ta{டயர்} mean? (\textquotedblleft{}A tyre\textquotedblright{}.) Say it once more.
\end{tcolorbox}
\section[\texorpdfstring{\ta{தீர்} (tīr)}{தீர் (tīr)}]{\ta{தீர்} (tīr) — to run out; to be used up}

\label{lesson:TA-C406-tir}



\begin{quote}
Before the new one: say the Tamil for petrol, then the Tamil for a tyre.
\end{quote}

\subsection*{You'll want to know: \ta{தீர்}}
\textbf{\ta{தீர்}} — \emph{tīr} — \textquotedblleft{}to run out; to be used up\textquotedblright{}.

\begin{quote}
\ta{நடுவழியில்} \ta{பெட்ரோல்} \ta{தீர்ந்தது}.
\end{quote}

\begin{quote}
\emph{naṭuvaḻiyil peṭrōl tīrntatu.}
\end{quote}

\begin{quote}
The petrol ran out halfway.
\end{quote}

Past (it …): \ta{தீர்ந்தது}, \emph{tīrntatu}.

\ta{தீர்வு}, a solution, comes from the same verb, which also means to be settled.

\begin{cousinweb}
\textbf{\ta{தீர்}} (\emph{tīr}) is native Dravidian (DEDR 3278), \textquotedblleft{}to come to an end, to be completed\textquotedblright{}.
\end{cousinweb}

\begin{grammarlens}[title={What you've built}]
One more word for getting around.
\end{grammarlens}

\subsection*{Your turn}
\begin{itemize}
\raggedright
  \item \emph{Say it:} \emph{tīr}
  \item \emph{Say it:} \emph{tīr}, once more
  \item \emph{Say it:} \emph{naṭuvaḻiyil peṭrōl tīrntatu.} — The petrol ran out halfway.
  \item \emph{Recall:} say the Tamil for petrol, then the Tamil for a tyre, then say \emph{tīr} again
\end{itemize}

\begin{tcolorbox}[breakable,colback=teal!4,colframe=teal!35!black,title={Before you move on}]
What does \ta{தீர்} mean? (\textquotedblleft{}To run out; to be used up\textquotedblright{}.) Say it once more.
\end{tcolorbox}
\section[\texorpdfstring{\ta{தேய்} (tēy)}{தேய் (tēy)}]{\ta{தேய்} (tēy) — to wear down; to rub}

\label{lesson:TA-C406-tey}



\begin{quote}
Before the new one: say the Tamil for a tyre, then the Tamil for to run out.
\end{quote}

\subsection*{You'll want to know: \ta{தேய்}}
\textbf{\ta{தேய்}} — \emph{tēy} — \textquotedblleft{}to wear down; to rub\textquotedblright{}.

\begin{quote}
\ta{இந்த} \ta{டயர்} \ta{நன்றாகத்} \ta{தேய்ந்துவிட்டது}.
\end{quote}

\begin{quote}
\emph{inta ṭayar naṉṟākat tēyntuviṭṭatu.}
\end{quote}

\begin{quote}
This tyre has worn right down.
\end{quote}

Past (it …): \ta{தேய்ந்தது}, \emph{tēyntatu}.

With an object it means to rub or scrub: \ta{பாத்திரம்} \ta{தேய்}, to scrub a pot.

\begin{cousinweb}
\textbf{\ta{தேய்}} (\emph{tēy}) is native Dravidian (DEDR 3458), \textquotedblleft{}to wear away by rubbing, to wane\textquotedblright{}.
\end{cousinweb}

\begin{grammarlens}[title={What you've built}]
One more word for getting around.
\end{grammarlens}

\subsection*{Your turn}
\begin{itemize}
\raggedright
  \item \emph{Say it:} \emph{tēy}
  \item \emph{Say it:} \emph{tēy}, once more
  \item \emph{Say it:} \emph{inta ṭayar naṉṟākat tēyntuviṭṭatu.} — This tyre has worn right down.
  \item \emph{Recall:} say the Tamil for a tyre, then the Tamil for to run out, then say \emph{tēy} again
\end{itemize}

\begin{tcolorbox}[breakable,colback=teal!4,colframe=teal!35!black,title={Before you move on}]
What does \ta{தேய்} mean? (\textquotedblleft{}To wear down; to rub\textquotedblright{}.) Say it once more.
\end{tcolorbox}
